\chapter{Viagra}
\index{Viagra}

\section*{Definition and Overview}
Viagra is the registered trademark name for sildenafil citrate, a potent and selective inhibitor of cyclic guanosine monophosphate (cGMP)-specific phosphodiesterase type 5 (PDE5). Initially synthesized in 1989 by Pfizer researchers Albert Wood and Peter Dunn as a cardiovascular treatment for hypertension and angina pectoris, sildenafil's therapeutic path was altered when clinical trials revealed its significant ability to induce penile erections. Recognizing its potential, Pfizer shifted its focus, leading to the landmark approval of Viagra by the United States Food and Drug Administration (FDA) in March 1998 as the first oral treatment for erectile dysfunction (ED).

Erectile dysfunction is defined as the persistent inability to achieve or maintain a penile erection sufficient for satisfactory sexual performance. Prior to the introduction of sildenafil, treatment options for ED were highly invasive and poorly tolerated, including intracavernosal injections, intraurethral suppositories, vacuum devices, or penile prostheses. Viagra revolutionized the management of ED, offering a non-invasive, highly effective, and generally safe oral option that could be taken on an as-needed basis.

In addition to its role in sexual medicine, sildenafil is also indicated under the brand name Revatio for the treatment of World Health Organization (WHO) Group 1 pulmonary arterial hypertension (PAH) to improve exercise capacity and delay clinical worsening. The introduction of sildenafil destigmatized erectile dysfunction, transforming public discussions about sexual health and relationships.

\section*{Epidemiology}
Erectile dysfunction is a highly prevalent medical condition that increases in frequency and severity with advancing age. The Massachusetts Male Aging Study (MMAS) established that approximately 52\% of men between the ages of 40 and 70 experience some degree of ED. This prevalence is strongly age-dependent: approximately 40\% of men are affected at age 40, rising to nearly 70\% by age 70. The proportion of men experiencing complete or severe ED increases from 5\% at age 40 to 15\% at age 70.

Globally, the epidemiological burden of ED is immense. Estimates suggest that the number of men suffering from ED worldwide was approximately 152 million in 1995, and this figure is projected to rise to 322 million by the year 2025, driven by an aging global population and the rising prevalence of comorbid metabolic and cardiovascular conditions. The prevalence of ED exceeds 75\% among men with diabetes mellitus, and it is significantly elevated in patients with hypertension, coronary artery disease, hyperlipidaemia, obesity, and chronic kidney disease.

For its secondary indication, pulmonary arterial hypertension (PAH), the epidemiology is distinct. PAH is a rare, orphan disease with an estimated incidence of 1.5 to 7.6 cases per million people per year, and a prevalence of 15 to 26 cases per million. Unlike ED, PAH disproportionately affects females, with a female-to-male ratio of approximately 3:1, peaking in the third to fifth decades of life.

\section*{Aetiology and Pathophysiology}
The physiology of a penile erection is a complex neurovascular event. Sexual stimulation triggers the release of nitric oxide (NO) from non-adrenergic, non-cholinergic (NANC) parasympathetic nerves and vascular endothelial cells within the corpus cavernosum of the penis. Nitric oxide diffuses into the smooth muscle cells of the corpus cavernosum and activates the enzyme soluble guanylate cyclase (sGC), which catalyzes the conversion of guanosine triphosphate (GTP) into cyclic guanosine monophosphate (cGMP).

Accumulated intracellular cGMP acts as a secondary messenger that activates cGMP-dependent protein kinase G (PKG), leading to a decrease in intracellular calcium levels and the relaxation of the trabecular smooth muscle of the corpus cavernosum and the smooth muscle of the helicine arteries. The subsequent vasodilation allows a rapid influx of arterial blood into the cavernous sinuses. The expanding erectile tissue compresses the subtunical venules against the rigid tunica albuginea, preventing venous outflow. This veno-occlusive mechanism traps blood within the corpus cavernosum, raising pressure and producing a rigid erection.

Sildenafil acts as a competitive, selective inhibitor of PDE5, the predominant enzyme responsible for the degradation of cGMP into inactive 5'-GMP in the corpus cavernosum. By inhibiting PDE5, sildenafil prevents the breakdown of cGMP, thereby enhancing the vasodilatory effects of nitric oxide. Sildenafil does not directly cause penile erections or increase libido; it only facilitates an erection in the presence of endogenous nitric oxide release, which requires active sexual stimulation.

The aetiology of erectile dysfunction is multifactorial and can be categorized into several primary groups:
\begin{itemize}
    \item \textbf{Vasculogenic Factors}: Atherosclerosis, peripheral arterial disease, and hypertension reduce blood flow to the penile arteries. Endothelial dysfunction, common in cardiovascular disease and diabetes, impairs the synthesis and release of nitric oxide.
    \item \textbf{Neurogenic Factors}: Damage to the autonomic pathways can occur due to spinal cord injuries, multiple sclerosis, Parkinson's disease, diabetic neuropathy, or pelvic surgeries such as radical prostatectomy.
    \item \textbf{Endocrinological Factors}: Hypogonadism (low testosterone) reduces nitric oxide synthase activity, while hyperprolactinaemia suppresses gonadotropin-releasing hormone, impairing erectile function.
    \item \textbf{Anatomical and Structural Factors}: Peyronie's disease causes fibrous plaques in the tunica albuginea, leading to painful penile curvature and structural ED.
    \item \textbf{Psychogenic Factors}: Performance anxiety, relationship discord, depression, general stress, and psychiatric conditions cause excessive sympathetic nervous activation, leading to vasoconstriction and preventing smooth muscle relaxation.
    \item \textbf{Pharmacological Factors}: Many medications can cause ED, including beta-blockers, thiazide diuretics, selective serotonin reuptake inhibitors (SSRIs), tricyclic antidepressants, antipsychotics, and anti-androgens.
    \item \textbf{Pulmonary Pathophysiology}: In PAH, chronic vasoconstriction, vascular remodeling, and smooth muscle hypertrophy in the pulmonary arteries restrict blood flow. Sildenafil inhibits PDE5 in the pulmonary vasculature, increasing cGMP and promoting pulmonary vasodilation, which reduces pulmonary arterial pressure and alleviates the workload on the right ventricle.
\end{itemize}

\section*{Clinical Features}
The clinical presentation of erectile dysfunction is characterized by the patient's subjective report of difficulty in sexual performance. Similarly, when sildenafil is prescribed for pulmonary arterial hypertension, a different set of clinical features is observed.

For patients presenting with erectile dysfunction, clinical features include:
\begin{itemize}
    \item \textbf{Inability to Achieve Erection}: The complete or partial inability to obtain a penile erection sufficient to initiate sexual intercourse.
    \item \textbf{Inability to Maintain Erection}: Achieving an erection initially, but losing rigidity before completion of the sexual act.
    \item \textbf{Reduced Penile Rigidity}: The erection is attained but lacks the necessary hardness to allow penetration.
    \item \textbf{Absence of Nocturnal or Morning Erections}: Healthy males typically experience involuntary erections during rapid eye movement (REM) sleep. The absence of these erections is a strong clinical indicator of an underlying organic (vascular or neurogenic) pathology rather than a psychogenic one.
    \item \textbf{Psychological Distress}: ED is frequently accompanied by low self-esteem, performance anxiety, relationship difficulties, and symptoms of clinical depression or anxiety.
    \item \textbf{Sudden vs. Gradual Onset}: A sudden onset of ED, particularly if situational (e.g., successful erection during masturbation but failure with a partner), points towards a psychogenic cause. A gradual, progressive decline in erectile function suggests an organic, vascular, or neurological progression.
\end{itemize}

For patients presenting with pulmonary arterial hypertension, the clinical features relate to impaired oxygenation and right-sided heart strain:
\begin{itemize}
    \item \textbf{Exertional Dyspnoea}: Shortness of breath during mild physical activity, which progresses to dyspnoea at rest as the disease advances.
    \item \textbf{Fatigue and Lethargy}: Generalized weakness due to reduced cardiac output and tissue hypoxia.
    \item \textbf{Syncope or Presyncope}: Fainting or feeling lightheaded, particularly during exertion, caused by the inability of the right ventricle to increase blood flow through the constricted pulmonary vascular bed.
    \item \textbf{Chest Pain}: Angina-like chest discomfort resulting from right ventricular ischemia due to hypertrophied right heart muscle working against high pressures.
    \item \textbf{Peripheral Oedema}: Swelling in the lower extremities, ankles, and feet, along with ascites (abdominal fluid accumulation), indicating progressive right-sided heart failure (cor pulmonale).
\end{itemize}

\section*{Differential Diagnosis}
When evaluating a patient presenting with erectile dysfunction, it is vital to differentiate the primary vascular or neurogenic causes of ED from other conditions that present with sexual dysfunction or physical limitations. Similarly, in the context of pulmonary hypertension, the underlying etiology must be carefully classified.

The differential diagnoses for erectile dysfunction include:
\begin{itemize}
    \item \textbf{Hypogonadism}: A deficiency in testosterone production that presents primarily with a profound loss of libido (sexual desire), fatigue, and loss of muscle mass. While ED is often present, the primary symptom is the lack of interest in sex, whereas patients with primary vasculogenic ED typically maintain a normal libido but cannot physically perform.
    \item \textbf{Hyperprolactinaemia}: Elevated serum prolactin, often caused by a pituitary prolactinoma or certain medications, which secondarily suppresses testosterone levels and impairs sexual function.
    \item \textbf{Premature Ejaculation}: A distinct form of male sexual dysfunction characterized by ejaculation occurring before the individual or partner wishes. Patients may confuse the rapid loss of erection post-ejaculation with primary erectile dysfunction.
    \item \textbf{Peyronie's Disease}: A condition marked by the development of fibrous scar tissue plaques under the skin of the penis, causing a significant, painful curvature during erection, which can physically impede penetration even if erectile rigidity is intact.
    \item \textbf{Isolated Psychogenic Erectile Dysfunction}: Erectile failure arising purely from psychological issues, relationship conflict, or performance anxiety in the absence of any vascular, neurogenic, or hormonal pathology.
    \item \textbf{Vascular Sentinel Effect of Coronary Artery Disease}: ED can be the initial presenting symptom of systemic atherosclerosis. Due to the small diameter of the cavernous arteries (1 to 2 mm) compared to coronary arteries (3 to 4 mm), arterial narrowing manifests earlier in the penis. Thus, ED must be differentiated from benign aging and investigated as a potential warning sign of impending myocardial infarction.
\end{itemize}

For pulmonary hypertension, sildenafil is only indicated for Group 1 PAH. Differential diagnoses include:
\begin{itemize}
    \item \textbf{Group 2 Pulmonary Hypertension}: Secondary to left heart disease (such as mitral valve disease or left ventricular systolic dysfunction). Treating these patients with sildenafil can be harmful, as it may cause pulmonary venous congestion.
    \item \textbf{Group 3 Pulmonary Hypertension}: Secondary to lung diseases or chronic hypoxia, such as chronic obstructive pulmonary disease (COPD) or interstitial lung disease (ILD).
    \item \textbf{Group 4 Pulmonary Hypertension}: Chronic thromboembolic pulmonary hypertension (CTEPH), which is treated primarily with surgical pulmonary endarterectomy rather than vasodilating medications.
\end{itemize}

\section*{When to Seek Medical Care}
While sildenafil is widely used and generally well-tolerated, there are several clinical scenarios and adverse drug reactions that represent acute medical emergencies. Patients must be educated on when to seek immediate medical intervention.

Seek immediate emergency medical attention if the following symptoms occur:
\begin{itemize}
    \item \textbf{Priapism}: An erection that lasts longer than 4 hours, which may or may not be painful. Priapism is a urological emergency. Ischaemic priapism leads to hypoxia and acidosis within the cavernous tissue, which can cause irreversible corporeal fibrosis and permanent, untreatable erectile dysfunction within 24 hours.
    \item \textbf{Sudden Vision Loss}: A sudden decrease or complete loss of vision in one or both eyes. This may be a sign of non-arteritic anterior ischaemic optic neuropathy (NAION), a condition where blood flow to the optic nerve is compromised. Patients must immediately discontinue sildenafil.
    \item \textbf{Sudden Hearing Loss}: A sudden reduction or loss of hearing, which may be accompanied by tinnitus (ringing in the ears) or vertigo.
    \item \textbf{Chest Pain or Cardiac Symptoms}: Any chest pain, pressure, tightness, or pain radiating to the jaw, neck, or arm during or after sexual activity. If emergency services are called, patients must immediately disclose if they have taken sildenafil within the past 24 hours. The administration of organic nitrates (e.g., glyceryl trinitrate) by emergency personnel can trigger profound, life-threatening hypotension.
    \item \textbf{Severe Hypersensitivity Reactions}: Signs of anaphylaxis, including facial swelling, urticaria, difficulty breathing, or severe dizziness.
\end{itemize}

For life-threatening symptoms, patients should immediately call their local emergency number, such as:
\begin{itemize}
    \item \textbf{local emergency services} in many countries.
    \item \textbf{911} in the United States and Canada.
    \item \textbf{999} in the United Kingdom.
    \item \textbf{112} in the European Union and other international regions.
\end{itemize}

\section*{Diagnosis and Investigations}
The diagnosis of erectile dysfunction is primarily clinical, established through a detailed medical, sexual, and psychosocial history, complemented by a physical examination and targeted laboratory investigations. The diagnostic workup is designed to identify the severity of the condition, screen for underlying cardiovascular and endocrine disease, and identify correctable risk factors.

The diagnostic process and investigations include:
\begin{itemize}
    \item \textbf{Clinical Questionnaire}: The International Index of Erectile Function (IIEF) is a validated, multi-dimensional self-reporting scale used to assess erectile function, orgasmic function, sexual desire, intercourse satisfaction, and overall satisfaction. The 5-item version (IIEF-5 or Sexual Health Inventory for Men - SHIM) is widely used in clinical practice to grade ED severity.
    \item \textbf{Physical Examination}: A focused examination includes measuring blood pressure and heart rate to assess cardiovascular status. A genital examination is performed to assess testicular volume, check for penile plaques (Peyronie's disease), and evaluate secondary sexual characteristics. Neurological assessment may include testing the bulbocavernosus reflex.
    \item \textbf{Endocrine Testing}: A fasting morning blood sample is drawn to measure total serum testosterone. If testosterone levels are low (typically less than 300 ng/dL or 10.4 nmol/L), the test should be repeated, and additional markers such as free testosterone, luteinizing hormone (LH), follicle-stimulating hormone (FSH), and prolactin should be measured to distinguish primary from secondary hypogonadism.
    \item \textbf{Metabolic Screening}: Fasting blood glucose or glycated haemoglobin (HbA1c) is measured to screen for undiagnosed diabetes mellitus. A fasting lipid profile (total cholesterol, HDL, LDL, and triglycerides) is obtained to evaluate cardiovascular risk.
    \item \textbf{Renal and Hepatic Function Tests}: Serum creatinine, urea, and liver enzymes are measured. Impaired renal or hepatic function alters sildenafil clearance, necessitating dose adjustments (e.g., initiating at 25 mg instead of 50 mg).
    \item \textbf{Penile Duplex Ultrasonography}: For patients with suspected vascular disease who do not respond to oral therapy, colour duplex Doppler ultrasound of the penile arteries can be performed after an intracavernosal injection of a vasoactive agent (e.g., alprostadil). This allows the measurement of peak systolic velocity (PSV) and end-diastolic velocity (EDV) to diagnose arterial insufficiency or veno-occlusive dysfunction.
    \item \textbf{Nocturnal Penile Tumescence (NPT) Monitoring}: Using a computerized device (such as the Rigiscan) worn overnight, this test records the frequency, duration, and rigidity of erections during sleep. The preservation of normal nocturnal erections confirms the integrity of the vascular and neurological pathways, indicating a primarily psychogenic etiology for the patient's daytime ED.
\end{itemize}

\section*{Management}
The management of erectile dysfunction is structured in a step-wise approach, addressing underlying cardiovascular and metabolic risk factors, providing psychological support, and prescribing appropriate pharmacotherapy or physical interventions.

The management options for erectile dysfunction include:
\begin{itemize}
    \item \textbf{Lifestyle Modification}: As a crucial first step, patients are encouraged to adopt healthy habits that improve endothelial function and cardiovascular health. These include smoking cessation, regular moderate-to-vigorous physical exercise, weight reduction in patients who are overweight or obese, and a Mediterranean-style diet.
    \item \textbf{First-Line Pharmacotherapy (PDE5 Inhibitors)}: Sildenafil (Viagra) remains the gold standard oral therapy. The typical starting dose is 50 mg, taken approximately 1 hour before anticipated sexual activity. Based on response and side effects, the dose may be increased to 100 mg or decreased to 25 mg. It should be taken on an empty stomach, as a high-fat meal can delay absorption and clinical onset. Other PDE5 inhibitors include tadalafil, vardenafil, and avanafil.
    \item \textbf{Psychosexual Counseling}: For patients with psychogenic ED or significant performance anxiety, referral to a qualified therapist or sexologist is highly beneficial. Cognitive behavioural therapy (CBT) can help break the cycle of anxiety and erectile failure, and partner-involved therapy improves relationship dynamics.
    \item \textbf{Vacuum Erection Devices (VED)}: A non-pharmacological option that involves placing a plastic cylinder over the penis and using a hand-operated or battery-powered pump to create a vacuum. This draws blood into the corpora cavernosa, producing an erection. A tension ring is then slipped onto the base of the penis to retain the blood; this ring must not be left in place for more than 30 minutes.
    \item \textbf{Intraurethral Alprostadil (MUSE)}: A tiny pellet of alprostadil (prostaglandin E1) is inserted directly into the urethra using a disposable applicator. The drug is absorbed through the urethral mucosa into the cavernous bodies, promoting vasodilation.
    \item \textbf{Intracavernosal Injection (ICI) Therapy}: Patients are trained to inject vasoactive medications directly into the lateral side of the corpus cavernosum using a fine needle. The most common formulations include alprostadil monotherapy, bi-mix (papaverine and phentolamine), or tri-mix (alprostadil, papaverine, and phentolamine). This therapy is highly effective and bypasses the need for intact neurological pathways.
    \item \textbf{Surgical Penile Prosthesis Implantation}: Reserved for patients who have failed or have contraindications to pharmacological treatments. An inflatable penile prosthesis (consisting of two cylinders placed in the corpora cavernosa, a fluid reservoir placed in the lower abdomen, and a pump placed in the scrotum) or a malleable semi-rigid rod prosthesis is surgically implanted. This option offers high satisfaction rates but is irreversible and carries a risk of infection or mechanical failure.
\end{itemize}

\section*{Complications}
The complications associated with sildenafil therapy can be divided into common, benign side effects resulting from systemic PDE5 inhibition, and rare, severe medical complications. Additionally, the consequences of untreated underlying conditions must be considered.

The complications and adverse effects include:
\begin{itemize}
    \item \textbf{Systemic Vasodilation Side Effects}: The most common side effects of sildenafil are mild to moderate and transient. These include headache (occurring in up to 16\% of patients), facial flushing (10\%), dyspepsia (7\%), nasal congestion (4\%), and dizziness (2\%). These are mediated by the relaxant effect of sildenafil on smooth muscle in other vascular beds.
    \item \textbf{Visual Disturbances (Cyanopsia)}: Sildenafil has a low but significant cross-reactivity with phosphodiesterase type 6 (PDE6), an enzyme found in the retina that is involved in phototransduction. This can lead to transient visual changes, most commonly a mild blue-green tinting of vision (cyanopsia), increased sensitivity to light, or blurred vision.
    \item \textbf{Severe Hypotension and Cardiovascular Collapse}: Sildenafil potentiates the hypotensive effects of organic nitrates (such as nitroglycerin, isosorbide mononitrate, and amyl nitrite) and guanylate cyclase stimulators (such as riociguat). Co-administration can cause an acute, catastrophic drop in blood pressure, leading to syncope, myocardial infarction, or stroke.
    \item \textbf{Non-Arteritic Anterior Ischaemic Optic Neuropathy (NAION)}: A rare but serious vascular condition of the eye characterized by sudden, painless unilateral vision loss, which can progress to permanent blindness. Sildenafil use has been associated with NAION in patients with pre-existing vascular risk factors (e.g., a crowded optic disc, hypertension, diabetes).
    \item \textbf{Ischaemic Priapism}: A prolonged erection lasting over 4 hours. If untreated, the blood within the corpora cavernosa becomes sturdily deoxygenated, leading to tissue necrosis, corporeal fibrosis, and severe, irreversible erectile dysfunction.
    \item \textbf{Ototoxicity}: Rare cases of sudden hearing loss, often unilateral and accompanied by vestibular symptoms such as vertigo or tinnitus, have been reported in temporal association with PDE5 inhibitor use.
    \item \textbf{Psychological Dependency}: Some patients without organic ED may develop a psychological dependency on sildenafil, relying on the drug to alleviate performance anxiety and experiencing panic or erectile failure if they attempt intercourse without it.
\end{itemize}

\section*{Prognosis and Outcomes}
The prognosis for men diagnosed with erectile dysfunction has improved dramatically with the introduction of sildenafil and other modern therapeutic modalities. Sildenafil is highly effective, with clinical trial data demonstrating that it improves erectile function in approximately 70\% to 80\% of treated men, compared to approximately 20\% in those taking a placebo.

However, the clinical response is heavily influenced by the patient's underlying aetiology:
\begin{itemize}
    \item \textbf{Psychogenic ED}: Has an excellent prognosis. With a combination of short-term sildenafil use and psychosexual counseling, many patients fully recover normal erectile function and can discontinue pharmacotherapy.
    \item \textbf{Vasculogenic ED}: Represents a chronic vascular condition. While sildenafil successfully manages the symptoms, it does not cure the underlying atherosclerosis. Efficacy may gradually decline over years if the underlying cardiovascular disease progresses.
    \item \textbf{Diabetic ED}: Patients with diabetes mellitus often have lower response rates to sildenafil (approximately 50\% to 60\%) due to a combination of severe endothelial dysfunction and advanced diabetic neuropathy. These patients may require higher doses (100 mg) or transition to combination therapies.
    \item \textbf{Post-Prostatectomy ED}: The prognosis depends on whether a nerve-sparing surgical technique was utilized. If bilateral nerves were spared, early penile rehabilitation protocols using sildenafil can lead to a recovery of spontaneous erectile function within 12 to 24 months. If non-nerve-sparing surgery was performed, oral sildenafil is rarely effective, and patients must rely on intracavernosal injections or penile implants.
\end{itemize}

Crucially, ED is now widely recognized as a sentinel marker for systemic vascular disease. The vascular diameter hypothesis outlines that since the penile arteries are small, they are obstructed by atherosclerotic plaque much earlier than the larger coronary or carotid arteries. Epidemiological studies show that the onset of ED precedes a major adverse cardiac event (such as myocardial infarction or stroke) by an average of 2 to 5 years. Therefore, the long-term prognosis of a patient presenting with ED depends on the comprehensive management of their systemic cardiovascular risk factors.

For pulmonary arterial hypertension (PAH), sildenafil (as Revatio) significantly improves functional capacity and haemodynamics. Although PAH remains a progressive, life-limiting illness, the integration of sildenafil into therapeutic regimens has helped increase the median survival time of patients from approximately 2.8 years in the early 1990s to over 5 to 7 years in modern cohorts, with some patients maintaining stable functional status for a decade or longer.

\section*{Prevention and Self-Care}
Preventing erectile dysfunction and optimizing the safety and efficacy of sildenafil therapy involves a combination of cardiovascular self-care, healthy lifestyle choices, and the responsible use of prescribed medications.

Key prevention and self-care strategies include:
\begin{itemize}
    \item \textbf{Dietary Management}: Adopting a diet high in phytonutrients, antioxidants, and unsaturated fats, such as the Mediterranean diet (emphasizing whole grains, vegetables, legumes, fruits, nuts, and olive oil, while limiting red meat and refined sugars). This diet preserves endothelial function and protects against vascular disease.
    \item \textbf{Regular Aerobic Activity}: Engaging in moderate to vigorous physical exercise (such as brisk walking, running, cycling, or swimming) for at least 150 minutes per week. Exercise enhances nitric oxide synthase activity and improves systemic perfusion.
    \item \textbf{Smoking Cessation}: Avoiding all tobacco products. Nicotine is a potent vasoconstrictor and direct endothelial toxin that severely impairs the vascular mechanisms necessary for erection.
    \item \textbf{Alcohol Moderation}: Limiting alcohol intake. While small amounts of alcohol can reduce anxiety, excessive consumption acts as a central nervous system depressant and can cause acute erectile failure.
    \item \textbf{Self-Monitoring of Metabolic Parameters}: Patients should undergo regular screening for blood pressure, blood glucose, HbA1c, and lipid profiles. Maintaining optimal control of these parameters prevents the microvascular and neuropathic complications that cause ED.
    \item \textbf{Stress and Mental Health Management}: Practicing stress-reduction techniques, such as mindfulness, meditation, or yoga. Addressing relationship conflicts openly and seeking couples therapy or sex therapy early can prevent the development of chronic psychogenic ED.
    \item \textbf{Safe Use of Sildenafil}: Patients must only obtain sildenafil via a legitimate prescription from a registered medical practitioner. Purchasing ``generic Viagra'' or herbal alternatives from unregulated online sources carries severe risks, as these counterfeits often contain incorrect dosages, toxic contaminants, or undeclared active pharmaceutical ingredients.
    \item \textbf{Proper Drug Administration}: Take sildenafil approximately 1 hour before sexual activity on an empty stomach or after a light meal. Avoid taking the medication with grapefruit or grapefruit juice, as grapefruit inhibits the cytochrome P450 3A4 (CYP3A4) enzyme, which metabolizes sildenafil, leading to elevated drug levels and an increased risk of side effects.
\end{itemize}

\section*{Sources and Further Reading}
\begin{itemize}
    \item World Health Organization (WHO) - Cardiovascular diseases: \url{https://www.who.int/health-topics/cardiovascular-diseases}
    \item American Urological Association (AUA) - Guideline on Erectile Dysfunction: \url{https://www.auanet.org/guidelines-and-quality/guidelines/erectile-dysfunction-guideline}
    \item Mayo Clinic - Erectile dysfunction: \url{https://www.mayoclinic.org/diseases-conditions/erectile-dysfunction/symptoms-causes/syc-20355776}
    \item National Institute of Diabetes and Digestive and Kidney Diseases (NIDDK) - Erectile Dysfunction: \url{https://www.niddk.nih.gov/health-information/urologic-diseases/erectile-dysfunction}
    \item European Association of Urology (EAU) - Sexual and Reproductive Health Guidelines: \url{https://uroweb.org/guidelines/sexual-and-reproductive-health}
    \item British Heart Foundation - Erectile dysfunction and heart disease: \url{https://www.bhf.org.uk/informationsupport/heart-matters-magazine/medical/erectile-dysfunction}
\end{itemize}